\newpage
\appendix
\onecolumn

\section{Example Questions Aligned with Bloom's Taxonomy}
\label{sec:appendix-questions}

\begin{table}[h!]
\centering
\small
\caption{Examples of questions aligned with Bloom's Taxonomy levels used in the study. The examples are drawn from visualisation \texttt{A}, with questions ranging from multiple-choice formats to open-ended tasks for the \textit{Create} level.}
\begin{tabular}{p{3.4cm}p{4.4cm}p{6.5cm}}
\toprule
\textbf{Level \& Definition} & \textbf{Example Question} & \textbf{Possible Answers} \\ 
\midrule
\textit{Identify}: Retrieving relevant knowledge (e.g., Recognising) &Which Empire had overseas colonies as early as the 1460s, as shown in the graph? & French; Spanish; \textbf{Portuguese}; British \\ 
\midrule
\textit{Understand}: Determining the meaning of instructional messages, including oral, written, and graphic communication (e.g., Explaining) & Which is the most probable explanation for the observed trend in the Spanish Empire line from 1800-1850? & Increased military conquests by Spain in Africa; \textbf{Wars of independence in Latin America}; Napoleonic Wars that led to the cession of colonies to France; The liberal principles in the 1812 Constitution were implemented throughout the Spanish Empire \\ 
\midrule
\textit{Apply}: Carrying out or using a procedure in a given situation. (e.g.,Implementing) & How would the graph change if the British Empire had captured half of Spain's colonies in 1700? Select all that apply. & \textbf{The distance between the French plot line and the British plot line would increase in 1700}; The distance between the British plot line and the Spanish plot line would decrease in 1700; The British plot line would reach its maximum in 1700; The Portuguese plot line would be greater than the Spanish plot line in 1700; The distance between the Portuguese plot line and the Spanish plot line would increase in 1700 \\ 
\midrule
\textit{Analyse}: Breaking material into its constituent parts and detecting how the parts relate to one another and to an overall structure or purpose. (e.g, Differentiating, Organizing)&According to the graph, how does the colonisation pattern of the French Empire compare to that of the British Empire? Select all that apply. & The British Empire always had more colonies than the French Empire; The British Empire did not lose colonies before 1900, while the French Empire did; \textbf{The British and French Empires started to decline simultaneously}; The British and French Empires reach their peak number of colonies in the same year; The British Empire gained overseas colonies years before the French Empire \\ 
\midrule
\textit{Evaluate}: Making judgments based on criteria and standards. (e.g., Critiquing) & What other pieces of information would you need to evaluate the validity of the following sentence: ``\textit{The Spanish Empire impacted negatively more people than the French Empire}''. Select only the relevant options. & The duration of colonialism under the French and Spanish Empires; The number of colonies over time for the British and French Empires; \textbf{The population size of the colonies under the French and Spanish Empires}; \textbf{Economic and cultural policies implemented by the French and Spanish Empires in their colonies}; The information on the graph is sufficient \\  \midrule
\textit{Create}: Putting elements together to form a novel, coherent whole or make an original product. (e.g., Generating, Planning, Producing) & Based on the visualisations, write a text on the relationship between the topics presented in the graphs. Your response should be 8 to 12 sentences long and should reference specific data points or trends from the graphs to support your ideas. & \textbf{[Open Ended Task]} Example of an extract of an answer:      \textit{The two figures can be linked by the historical trends.  Figure 1 shaped the data shown in Figure 2. For example, Figure 1 shows the rise and fall of Spanish colonialism. We see likely downstream consequences of this reflected in Figure 2, with certain countries having higher rates of uncontrolled territory, as with Mexico, Peru, Colombia, Venezuela.} \\
\bottomrule
\end{tabular}
\label{tab:question_table}
\end{table}

\newpage
\section{Create Task Labelling Rubric}
\label{sec:appendix-create-rubric}

\begin{table}[h]
\centering
\caption{Rubric used to categorise sentences from participants' responses in the \textit{Create} task, based on the cognitive levels of Bloom's Taxonomy. }
\small
\begin{tabular}{p{1.5cm}p{5cm}p{8cm}}
\toprule
\textbf{Level} & \textbf{Description} & \textbf{Example} \\ \midrule
\textit{Identify} & Factual information from a graph without interpretation. Reference to values and information in one graph. & For example, the Spanish started colonizing the Americas in the 1400s and the British and French in the 1600s as can be seen from the first graph. \\ \midrule
\textit{Understand} & Explaining only one graph and interpreting the values. & The factors that affect GDP of a country including civil wars, politics, government policies, resources, investments and local and international trade. \\ \midrule
\textit{Apply} & Applying information from a graph to a new situation or hypothetical scenario about future trends. & For Syria, the civil war is still ongoing, so it is likely their state capacity is still low in 2022. \\ \midrule
\textit{Analysis} & Explanation linking both topics (reference to two graphs). Identifying potential relationships between the two topics. & The length of time that former colonies have been independent correlates with how much control those countries have over their territory. \\ \midrule
\textit{Evaluate} & Critical assessment of information and underlying assumptions in the graphs. Highlighting discrepancies, outliers or missing elements that affect interpretation. & More information is needed to validate the relationship as the first graph provides information until 2010 and the second graph provides information of 2022. \\ \midrule
\textit{Create} & Extending or extrapolating beyond the information provided in the graph to draw new conclusions, or propose alternative interpretations in the form of a new argument. & At the same time, rather than stability, countries with high rates of controlled land (that were post-colonial), such as South Africa, might not necessarily be a positive thing as colonialism imposed cultural, social, political norms, perhaps they had less ability for self-determination (versus other countries that had more ability to do so, even if it means higher rates of uncontrolled territory). \\ \midrule
\textit{Other} & Information not directly related to the graph to introduce another sentence (for example: topic sentences). Sentences of previous knowledge. & It's interesting to compare these visualisations between these two countries. \\ \bottomrule
\end{tabular}
\label{tab:levels}
\end{table}